\chapter{Stories Are Older Than Writing}
\section{A Song You Never Memorised}
Pick up something: your headphones.

Put them on and play a song you often hear. Two seconds into the introduction, your mind supplies the next line. At the chorus, you can sing every word, although you never memorised it like a school text. Nobody corrected your spelling or tested your recall. Yet you know it firmly enough that it sometimes loops inside your head and refuses to leave.

Compare that with a twenty-character, four-line poem in Chinese class. You spent a whole morning study session learning it, yet reversed two lines in the examination.

Your memory is not applying double standards. Here is the chapter's first secret: human minds naturally remember certain things, especially rhythm, melody, repetition, and material hung on a storyline. Lyrics possess all these hooks; the school passage does not. You think you remembered without memorising, but the songwriter, arranger, and thousands of years of singing traditions installed the hooks for you.

Scale that experience up to humanity. Writing arrived late. The earliest generally recognised mature systems, Sumerian cuneiform and ancient Egyptian hieroglyphs, appeared a little over five thousand years ago. Human speech goes back at least tens of thousands of years. For an immense span, then, all knowledge, history, law, myth, genealogy, and ancestral deeds had no written words. Everything lived in heads, on tongues, and in songs.

How did hundreds of thousands of generations without paper, pens, or hard drives remember everything? How did epics, myths, legends, and songs cross millennia and survive into the book you hold?

\section{Singing in a Coffeehouse}
Come into another scene.

The time is the 1930s; the place, a coffeehouse in a mountain town in Yugoslavia, around present-day Bosnia. Smoke fills the room. Men sit drinking strong coffee. In a corner stand several recording devices brought from America, still uncommon then, with long recording wires wound around aluminium discs.

Their owner is an American scholar, Milman Parry, who has travelled far to answer a problem unresolved for more than two thousand years. We will examine it shortly. His answer sits in a chair in the coffeehouse's centre: a singer holding an instrument.

The instrument is a gusle. Its single string sounds monotonous, hardly melodic, more like a beat for the voice. Most singers come from ordinary herding or farming families; many cannot read. One of the most famous is Avdo Međedović, whom Parry's student Albert Lord also recorded later.

What can Avdo do? Sing epics. An epic can run to more than ten thousand lines. He sings from evening deep into night, resumes the next day, and takes several days to finish. His subjects are wars, heroes, weddings, and raids from centuries earlier: how the warrior Smailagić Meho won his bride, how a hero crossed the Drina to seek revenge. Listeners drink coffee, cheer at exciting moments, or doze. The singer does not mind; he extends a passage or simply skips it.

Lord asks everyone's question: how did you memorise such a long poem?

Avdo's answer astonished Western scholarship for decades. He had not memorised it. He had learned no ``original text'', because what he sang had no fixed original. He was \textbf{remaking it as he sang}.

To test this, scholars asked him to sing the same epic again. He did. The story and heroes remained, but a word-by-word comparison showed large differences: longer here, shorter there, passages in different positions. More astonishingly, Avdo saw no problem. In his world the performances were obviously ``the same song'', just as tomato-and-egg stir-fry cooked today and tomorrow is the same dish, although no two tomato slices are identical.

Remember this coffeehouse. The rest of the chapter answers the questions it throws at us.

\section[Remembering Ten Thousand Lines]{How Can a Nonreader Remember Ten Thousand Lines?}
Lay out the conflict without rushing to an answer.

Two intuitions clash when we encounter a singer like Avdo.

The first says this is merely an exceptional memory, nothing remarkable. An illiterate person has no skill beyond rote learning. Moreover, oral transmission is unreliable: anyone who has played the telephone game knows a sentence becomes unrecognisable after ten people pass it along. Before writing, cultural memory must therefore have been crude, distorted, and liable to break at any moment. Humanity only truly began keeping records when writing appeared.

The second says no. Singing for days, producing ten thousand metrically regular lines with coherent structure and vivid characters, cannot be explained by ``good memory''. Even our best memorisers struggle with a passage of a few thousand words, let alone ten thousand lines that can be changed in performance. Some unfamiliar technology must be at work.

Behind these intuitions lie three genuine questions, not merely respect or contempt for other people:

\textbf{First: without writing, what does ``remembering'' mean?} Our generation understands memory as reproducing a fixed original: reciting schoolwork means matching the textbook. But what if no original exists? What if remembering a song means remembering the ability to \textbf{make it again whenever needed}?

\textbf{Second: when oral material changes as it travels, is that a defect or its way of working?} We use the telephone game to show oral unreliability, but that game transmits an ``original sentence''. A singer transmits a story. If it changes shape, is it still the same story?

\textbf{Third: what are myths and legends such as Nüwa repairing the sky, the Great Flood, or the woman whose tears collapsed the Great Wall?} ``Fake news'' invented by scientifically ignorant ancestors, or something else?

Before continuing, take a position. If writing vanished tomorrow, how far would knowledge retreat: back to a primitive society, or less catastrophically than we imagine?

\section[Writing, Speech, and a Weeping Woman]{Writing, Speech, and the Woman Who Wept Down the Wall}
Give both positions their fullest hearing.

\textbf{First, the strongest case for writing.} It deserves care precisely because writing's superiority to speech seems so obvious today that nobody thinks it needs an argument.

The case has at least three layers. First, writing remembers accurately. Oral transmission depends on minds that forget, err, and embellish. A written contract, ledger, prescription, or experimental dataset retains every word after a thousand years. Without this precision, modern science, law, and medicine would be impossible. Second, writing travels far. Spoken messages reach only listening ears; letters cross oceans, and books converse with people dead for two millennia. We read Confucius today because of bamboo slips and paper, not an unbroken chain of voices. Third, writing accumulates. Every generation need not begin again. Newton's famous thought was that he saw farther by standing on giants' shoulders, and libraries built those shoulders.

These reasons are too strong to pretend away. Acknowledging writing's power is our starting point, not our conclusion.

\textbf{Now hear an opposing witness, himself a giant of the written world.} In the \emph{Phaedrus}, Plato has Socrates tell a story. An Egyptian god invents writing and offers it to the king, claiming it will improve people's intelligence and memory. The king refuses: on the contrary, people will depend on external signs instead of remembering inwardly. Writing brings not memory but forgetfulness, not wisdom but its appearance.

Notice the strangeness. Plato wrote books, and we know this story by reading them. A man remembered through writing uses his teacher's voice to warn against writing. The argument raged even in writing's own homeland.

\textbf{Hear a third voice: ordinary people.} Consider the Chinese legend of Meng Jiangnü weeping at the Great Wall. You may know it. The First Emperor conscripts Wan Xiliang to build the Wall. His wife Meng Jiangnü travels a great distance with winter clothes, only to find him dead, his bones buried inside the wall. Her earth-shaking grief brings part of it down, revealing his remains.

Is this official history? No. But neither is it ``fake news'' dreamed up in a moment. Modern scholar Gu Jiegang spent years tracing its transformations. Its earliest seed was a brief account in the \emph{Zuo Commentary}: the Qi general Qi Liang died in battle during the Spring and Autumn period, and his wife refused the Duke of Qi's roadside condolences as ritually improper. There was no collapsing wall or First Emperor; even the date differed by centuries. Han accounts added weeping for the husband and a wall's collapse. Later generations repeatedly relocated the story: from Spring and Autumn to Qin, from Qi Liang's wife to Meng Jiangnü, with the husband buried in the Great Wall and criticism directed at the First Emperor.

See what happened? Over two thousand years of telling, each generation rewrote the story into what it needed. Other legends follow similar paths. The White Snake changes from a frightening demon in Tang anecdotal writing into today's passionately loving Lady White. The Cowherd and Weaver Girl grow from two star names into a full tale of separation and reunion. Almost every familiar story has a genealogy of changing forms. After Qin and Han, ordinary people dared not criticise emperors in memorials, but legends could let a woman's grief overturn tyranny's walls. \textbf{Legends are not counterfeit history; they are the only place ordinary people can write history.} Official histories record rulers' deeds; legends record ordinary people's feelings. Both matter. Without either, history is deaf in one ear.

\textbf{Finally, a fourth voice on myth.} Why do almost all cultures tell flood stories? China has Nüwa repairing heaven and Yu controlling floods; Greek myth has Deucalion and his wife escaping Zeus's flood by boat; western Asia's \emph{Epic of Gilgamesh} has an almost identical episode. Were ancient people conspiring to lie? No. Myth does not principally answer what happened on a particular date. It asks where the world came from, why disaster strikes, how humans relate to gods and nature, and who we are. Myth is a community's narrative for \textbf{organising meaning}. It weaves scattered experience, fear, and hope into a story with a beginning and end, telling a group where it came from and where it should go. Judging myth simply as true or false is like weighing a song to measure its length. The myth has not answered wrongly; the tool is wrong.

\section{Thinking Bridge: Hooks for Memory}
Here is a portable tool answering our first question: how does someone who cannot read remember ten thousand lines?

It is simple: \textbf{when you encounter something orally transmitted---a song, ballad, epic, or jingle---do not first ask its meaning. Look for three components: repetition, rhythm, and fixed phrases.}

First, \textbf{repetition}. Why can you not forget an earworm? Its chorus returns again and again. Repetition is not verbosity but memory's rivet. For listeners, a sentence disappears once spoken. They cannot turn back a page to inspect it. Oral cultures therefore repeat important things twice, three times, ten times, each repetition giving the audience another chance to catch them.

Next, \textbf{rhythm}. Reciting Li Bai's line about moonlight before the bed feels smoother than memorising prose because rhythm and rhyme lay tracks for sound, allowing the mind to slide along. Oral traditions worldwide know this: poems rhyme, songs keep time, and epics follow instruments. The gusle's monotonous string does exactly that. It produces not music but a beat, a track for the singer's voice.

Finally, \textbf{fixed phrases}, what scholars call ``formulas''. Since the term is new, consider familiar examples: a livestreamer's ``Friends, put up the link!'', the birthday song's slot for a person's name, or a storyteller's opening, ``The empire, long divided, must unite; long united, must divide.'' These are reusable, partly finished sentences, like Lego bricks. An oral singer's mind holds not rigidly memorised poems but thousands of such bricks and story frameworks: someone marries, seeks revenge, or crosses a river. In performance, the singer fits bricks into slots along the beat, short phrases for quick passages, long ones for slow passages, assembling on site without error.

Avdo's secret, then, is not extraordinary memory: \textbf{he does not memorise poems but the machine that makes poems}. Repetition, rhythm, and formulas are its blueprint.

This machine is everywhere. Have you memorised the twenty-four solar terms? Few learn the table, but many Chinese children hear the song whose opening syllables condense the spring and summer terms: spring, rain, awakening, equinox, clarity, grain; summer, fullness, awns, solstice, and the heats. Farmers need not open a calendar; start singing and the terms line up. Military-training call-and-response works similarly: ``One, two, three, four, five; we've waited such a long time!'' Numbers establish rhythm, and the second half can be replaced. A whole company joins without rehearsal. These use the same ancient technology as the \emph{Book of Songs} and Homer. Next time a song refuses to leave your head, take it apart: where are its hooks?

\section{The Secret in a Patch of Plantain}
Read a short primary text and dismantle it with the new tool.

The \emph{Book of Songs}, China's earliest poetry anthology, collects songs from roughly three thousand to twenty-five hundred years ago. They were originally \textbf{sung}: fossils of oral tradition fixed onto bamboo slips. Here is the complete ``Plantain'' from its ``Airs of Zhou and the South'':

\begin{quote}
Gather, gather plantain; come, let us gather it. Gather, gather plantain; come, let us have it. Gather, gather plantain; come, let us pick it. Gather, gather plantain; come, let us strip it. Gather, gather plantain; come, let us hold it in our skirts. Gather, gather plantain; come, let us tuck our skirts to carry it.
\end{quote}
Fuyi, pronounced fú yǐ, is plantain, a wild vegetable and medicinal herb. What ``happens''? Almost nothing. Women gather plantain, collecting more and more. The poem has three stanzas with identical sentence patterns, changing only two verbs per stanza: gather, have, pick, strip off by handfuls, hold in the garment's front, and tuck the garment into the belt to carry more.

Scholars call this ``repeated stanzas and refrains'': the same stanza returns with only a few words changed. It is oral tradition's fingerprint on paper. Before bamboo slips recorded them, generations sang these songs in fields and ceremonial halls. Their repetitive structure was designed for singing while working and learning at once.

Apply our tool. Repetition? The entire poem embodies it. Rhythm? Four Chinese characters per phrase, paired like a work chant's beat. Formula? ``Gather, gather plantain; come, let us X it'' is a standard slot. Change the action; keep the frame.

Read it twice more and this poem in which ``nothing happens'' becomes moving. Gather, have, pick, strip, scoop into clothing: the work gets faster, the crop larger. Repetition holds the workers' growing enjoyment. You can almost see people singing along the field bank as they work. There is nothing ``inferior'' about this technique. Minimal variation creates a remarkably full sense of presence.

Consider a more familiar example, Mulan preparing to depart in the \emph{Ballad of Mulan}:

\begin{quote}
At the eastern market, buy a fine horse; at the western market, a saddle and pad; at the southern market, a bridle; at the northern market, a long whip.
\end{quote}
Who actually visits four markets to buy one item at each? This is not a shopping list but formulaic elaboration. Four directional patterns sing the bustle and solemnity of departure. As a factual record it is absurd; as a song it is exactly right. These passages show that \textbf{oral literature has its own grammar. Judging it by written literature's standards is like refereeing basketball with football rules.}

\section{Who Was Homer? Three Answers}
Return to Parry's original puzzle, one of Western scholarship's most famous unresolved cases: the ``Homeric question''.

First, the factual layer: what happened? The \emph{Iliad} and \emph{Odyssey} are ancient Greek epics, each exceeding ten thousand lines, concerning the Trojan War and Odysseus's journey home. They took written form around the eighth century BCE and are attributed to Homer, a legendary blind bard. These are the parts supported by evidence. How Homer produced these works is an explanatory question. More than a century of argument has offered three main answers:

\textbf{Explanation one: individual genius.} Homer really existed, a great author like Dante or Shakespeare, independently composing or finally editing the epics. The argument is their astonishing structural unity and consistent characters. The \emph{Iliad}'s thousands of lines revolve tightly around Achilles's anger. Such artistic wholeness hardly resembles a group production: a committee cannot paint the \emph{Mona Lisa}. Its weakness is the unusually dense formulas. Dawn is always rosy-fingered, the sea wine-dark, and Athena almost always arrives with a fixed epithet. No individual author concerned with originality would write this way.

\textbf{Explanation two: accumulated tradition.} There was no individual Homer. The name is a collective signature of countless bards over centuries. Oral singing built the epics layer upon layer, like a coral reef rather than a sculpture. The evidence is strong: formulas are precisely oral composition's components, as we will see, and the language mixes dialect layers from different regions and periods, exposing age like geological strata. Its weakness is artistic unity. Why does something countless people casually added to and altered over centuries read like one mind's work?

\textbf{Explanation three: oral growth, written fixation.} The epics grew orally for centuries, absorbing countless singers' craftsmanship. Then, after writing reached Greece, one exceptionally gifted poet or a small group recorded, edited, and fixed them. This combines both sides: formulas come from oral tradition, unity from final editors. Most scholars today accept some version of this account. Its weakness is the scant evidence about that moment of fixation: who, where, and how? It remains a reasonable inference more than a proven fact.

Keep four things separate. The poems' existence and approximate eighth-century BCE written form are \textbf{facts}. The three accounts compete over \textbf{why}. Ancient Greeks' widespread belief in Homer as a real blind poet is \textbf{how people then understood it}. Whether we credit ``genius'' or ``tradition'' is \textbf{how we evaluate it}. The century-long argument matters not because an answer is imminent but because each generation sharpens the explanations with new evidence. This was the question Parry carried into our Yugoslav coffeehouse.

\section[Further Challenge: Every Performance Is Original]{Further Challenge: Every Performance Is an Original}
Parry found an answer in the coffeehouse. His and Lord's work became known as ``oral-formulaic theory'', this chapter's weightiest theory. Take time with it.

Parry first discovered in his study that Homeric formulas were not decorative language but parts of a metrical machine. Each Greek epic line has fixed rhythmic slots, and tradition supplies phrases of different lengths for almost every common character and scene. Short phrases fill short slots; long ones fill long slots. Singers need not ``think up words''. Familiar hands reach into the parts store and select by slot, assembling while singing. ``Rosy-fingered dawn'' is the standard component for dawn at a particular slot length.

Lord then verified this in Yugoslavia. Nonreading singers such as Avdo ran the same machine. They had no concept of an ``original text''. They knew story frameworks, thematic scenes such as assemblies, departures, banquets, and weddings, each with an established method of performance, and a stock of formulas. Each performance was composition on the spot.

The theory reaches a disruptive conclusion: \textbf{in oral tradition, no ``original work'' exists.} Written literature ranks the original above copies, translations, and adaptations. Yet Avdo can sing the same epic twice with substantially different words, and nobody calls one genuine and the other pirated. Every competent performance fully realises the song. Oral literature is not a book but a craft. It exists not in a text but in singers, as cooking exists in cooks. To collect it, you would have to collect a person.

This also resolves our second question: is continual change a defect or a working method? \textbf{Change is its working method.} Meng Jiangnü's journey from Spring and Autumn to Qin is not telephone-game distortion. Every generation composes afresh, adjusting the story to its own time. Oral stories are living things, and living things change.

Be careful, however, of two tempting misjudgements. Here are two counterexamples.

\textbf{Counterexample one: oral transmission need not distort.} India's \emph{Rigveda}, composed over three thousand years ago, depended entirely on oral transmission for a long period. Its recitation tradition trains rigorously and repeatedly checks every phonetic detail, preserving language with astonishing precision. Linguists can read forms in today's transmitted versions older than those in many surviving manuscript traditions. Oral transmission need not be crude. When precision is itself the goal, human minds can achieve hard-drive-like feats. Conversely, writing is no guarantee. The Grimm brothers' fairy tales appeared in print, yet from the 1812 first edition onwards they revised them repeatedly, removing excessively bloody episodes and changing abusive ``mothers'' into ``stepmothers''. \textbf{Written texts also change in transmission. It is not speech or writing that changes stories, but the act of telling stories to people.}

\textbf{Counterexample two: ``every performance is original'' does not mean anything goes.} Oral-formulaic theory describes a craft, not a licence to fabricate carelessly. Avdo's performances differ but obey strict traditional constraints. The framework cannot be scrambled, the metre cannot break, and knowledgeable listeners boo poor singing. Freedom within tradition is another lesson of oral literature.

Finally, the theory allows a fairer judgement between speech and writing. After fieldwork in the western Pacific, anthropologist Bronisław Malinowski argued, in essence, that myth is not ancient scientific speculation but a social ``charter''. It gives existing arrangements origins and reasons. Why does this family own that land? Our mythical ancestor landed there. Why perform this ritual that way? The song says so. In oral societies, stories unite archive, legal code, and constitution.

Oral and written literature are therefore not lower and higher stages but \textbf{two different technologies with different strengths}. Writing excels at precision, accumulation, travel across distances and millennia, and complex argument. Speech excels at presence, participation, emotional resonance, readily available memory, and stitching a community together.

Living evidence abounds. China's three great epics, Tibetan \emph{Gesar}, Mongolian \emph{Jangar}, and Kyrgyz \emph{Manas}, began as oral works. \emph{Gesar} is called a ``living epic''. Its total length far exceeds the \emph{Iliad} and \emph{Odyssey} together, and singers on the Tibetan Plateau still perform instalment after instalment. Some with little literacy can sing dozens. In West Africa, hereditary singer-historians called griots recite royal genealogies and perform \emph{Sunjata}, the founding epic of the thirteenth-century Mali Empire, at weddings and celebrations. That empire left no written national history; it entrusted history to song. In Australia, Aboriginal songlines weave geography, water sources, prohibitions, and ancestral deeds into songs. Singing one amounts to travelling a route and reviewing essential survival knowledge. None of these traditions regards itself as an unfinished product waiting to become a book. Each is a complete, functioning instrument of civilisation.

Many things writing cannot replace are among the most important. Otherwise, why do we still sing, tell jokes, hear podcasts, and watch short videos in an age of near-universal literacy?

\section{Memes, Rap, and a New Oral Age}
Indeed, scholars observe a large-scale return of oral culture.

Communication scholar Walter Ong called this ``secondary orality'': radio, television, and the internet foster a new oral character built on writing, yet transmitted, like ancient songs, through sound, immediacy, repetition, and collective participation. Everything in this chapter is alive today:

\textbf{Formulas survive.} Memes and reaction images are contemporary formulas: fixed, ready-to-use patterns instantly understood by insiders. Once ``Listen as I thank you'' becomes a formula, the remainder need not be spoken. Livestream calls of ``Friends!'' and ``Put up the link!'' are commercial society's equivalents of ``The empire, long divided, must unite''.

\textbf{On-the-spot assembly survives.} Freestyle rap battles fill words into a beat in real time: Avdo's craft with a microphone. Short-video choruses and ``Plantain'' use the same memory rivets.

\textbf{``Every performance is original'' survives.} A meme acquires three versions in three days, and nobody asks which social-media post was the original. Participants are audience and singers alike, each a link in tradition. Change one word while sharing and you join a relay thousands of years old. Voice messages to friends and bedtime podcasts are everyday secondary orality too. More than five thousand years after writing's invention, people still love accompanying one another with voices, whose tone, pauses, and warmth writing cannot contain.

\textbf{Legends survive, including their danger.} Urban legends in group chats and short videos, allegedly witnessed by ``a friend of a friend'', localise details in each city and mutate with anxieties. Their structure matches Meng Jiangnü's. This reminds us that a legend's ability to change is neutral. It can weep down tyranny's walls or drown an innocent person. Ancient craft wielded by new technology makes both memes and rumours. Knowing how stories mutate brings responsibility. Before forwarding, ask: is this story organising meaning for me, or manufacturing hatred for someone?

\section{For You: Whom Will You Trust with Memory?}
Return to the coffeehouse.

Avdo died many years ago. Fortunately, Lord's recording wires preserved his voice, and archives still let you hear that single-stringed instrument. But notice the sadness: the recording keeps \textbf{that one} performance. The ability to remake it whenever needed vanished with the singer. Writing and recordings preserve specimens, not living things.

Here is the final question. To transmit something to people a thousand years from now, you must choose one of two methods.

Write it down. Writing will lie silently in ruins, unafraid of floods or forgetting, patiently awaiting someone able to read it. But if nobody recognises the script after a millennium, it is only a pattern.

Or teach a group to sing it. Song lives in people, without electricity or libraries. Where people remain, the song remains, crossing mountains and seas and growing and changing by itself. Yet it depends on a generational relay. Wherever the chain breaks, the song dies.

Whom will you trust with memory?

Before answering, weigh writing's precision and silence against speech's vitality and fragility. Voices carried the \emph{Rigveda} across three thousand years, while an inscribed stone may become unreadable after two hundred. Homer's epics were born orally and saved in writing; King Gesar's epic is still sung into being line by line. Some memories seem to need both technologies to survive.

There is no standard answer, but give your reasons. The question concerns more than the past. Every day you vote: where do you write what deserves remembering, and to whom do you sing it?
